\section{Modal Participation Factors}

\subsection{Mass Participation Factors}

In addition to visualising the mode shapes, the mass participation factor table provided by ANSYS has been used. This table makes it possible to identify which modes have greater dynamic participation in each global direction of the model.

\begin{table}[H]
\centering
\caption{Modal participation factors in the principal directions.}
\label{tab:modal_participation}
\begin{tabular}{cccccc}
\toprule
Mode & $f$ [\si{\hertz}] & $M_x$ [\%] & $M_y$ [\%] & $M_z$ [\%] & Dominant direction \\
\midrule
1 & & & & & \\
2 & & & & & \\
3 & & & & & \\
4 & & & & & \\
5 & & & & & \\
6 & & & & & \\
\bottomrule
\end{tabular}
\end{table}

\subsection{Identification of the Most Critical Modes}

% Justify which modes are the most critical or relevant for the geometry,
% considering jointly:
% - The value of the natural frequency
% - The dominant direction of the mode
% - The modal participation factor
% - The proximity to possible excitation frequencies
% - The areas of largest deformation
% - The possible occurrence of critical displacements or stresses